\section{Kaggle implementation: K-means versus DBSCAN}
\label{sec:kaggle-mall-clustering}

\begin{kaggleproject}{Mall Customers --- forced partitions versus density-based noise}
\kagglesource{https://www.kaggle.com/datasets/vjchoudhary7/customer-segmentation-tutorial-in-python}{Mall Customer Segmentation Data}
\textbf{Question.} How does a method that assigns every row to one of $k$ clusters differ from a density method that can mark rows as noise?
\end{kaggleproject}

The program compares K-means with DBSCAN on the same standardized representation, reports silhouette evidence where meaningful, and summarizes the K-means segments in the original units.

\textbf{Before running.} Predict the qualitative pattern you expect from the experiment before you look at the output or plot.

\begin{labmode}
Stage 4B: choose which diagnostics and original-unit summaries you need before consulting the complete comparison script.
\end{labmode}

\lstinputlisting[style=mlpython]{all_experiments/lab15-mall_clustering.py}


\begin{tryit}
Vary DBSCAN's \texttt{eps} while keeping \texttt{min\_samples} fixed. Track the number of clusters and noise points. Explain why tuning DBSCAN only to maximize silhouette can still produce a segmentation that is operationally useless.
\end{tryit}

\begin{labdeliverable}
Submit one notebook comparing K-means and DBSCAN on the same standardized representation. Include cluster/noise counts, silhouette evidence where valid, an original-unit K-means profile table, and an \texttt{eps} sensitivity table showing how DBSCAN's cluster/noise structure changes. End with one paragraph on which result is more operationally interpretable and why.
\end{labdeliverable}
